% ── SECTION 1: INTRODUCTION (~500 words) ──────────────────────────────────────
% Cover: the interpretive gap, why a formal relational framework is needed,
%        preview of RCT and U_RA gate.

\section{Introduction}
\label{sec:intro}

% TODO: Expand to ~500 words.
% Key points to make:
%   1. Quantum mechanics is empirically unmatched but interpretively contested.
%   2. Bell's theorem (1964) + Zeilinger Nobel (2022) definitively rule out
%      local hidden variables — quantum properties are not intrinsic.
%   3. The dominant interpretations (Copenhagen, many-worlds, QBism) defer or
%      dissolve the ontological question rather than answering it.
%   4. Rovelli's relational QM (1996) takes the relational structure seriously
%      but has not been formalized as a mathematical framework with testable
%      predictions.
%   5. This paper provides that formalization: the RCT + Affinity Tensor α
%      + the U_RA gate.
%   6. Companion paper (Pub 2, philosophy) provides the full ontological argument.

The experimental confirmation of entanglement nonlocality, culminating in the
2022 Nobel Prize in Physics awarded to Aspect, Clauser, and
Zeilinger~\cite{Zeilinger2023Nobel,Aspect1982,CHSH1969}, has definitively
established that quantum mechanics cannot be completed by local hidden
variables~\cite{Bell1964}.
Quantum properties are not intrinsic to isolated systems; they emerge through
relational interaction~\cite{Rovelli1996,Einstein1935EPR}.

Despite this empirical clarity, quantum mechanics lacks a formal ontological
framework that makes the relational structure mathematically precise and
derives testable predictions from it.
Relational quantum mechanics~\cite{Rovelli1996,Rovelli2021} correctly identifies
the relational character of quantum states but has not been formalized into
a mathematical structure linking coherence, entanglement, and decoherence
dynamics under a unified relational concept.

We address this gap by introducing the \emph{Relational Coherence Theorem}
(\RCT) and the associated \emph{Affinity Tensor}~$\alpha(S,E)$.
The \RCT~formalizes the claim that relational affinity~---~the degree of
coherent coupling between a system and its environment~---~is the fundamental
quantity governing both the emergence of definite properties and the transition
from quantum to classical behavior under decoherence~\cite{Zurek2003}.

From the \RCT~we derive the \emph{Relational Affinity Gate}~$\URA(\alpha)$,
a continuously parameterized two-qubit gate whose entanglement capability is
tuned by the Affinity Tensor value.
Section~\ref{sec:background} reviews the required background.
Section~\ref{sec:rct} states and proves the \RCT.
Section~\ref{sec:gate} characterizes~$\URA$ in the Weyl chamber.
Section~\ref{sec:sim} presents numerical simulation results.
Sections~\ref{sec:disc}~and~\ref{sec:conc} discuss implications and conclude.
